\subsection{问题一：单日计划购电策略}

\subsubsection{问题分析}

问题一只给出一天的数据：电价曲线按题面要求逐日相同，而小区负载与光伏出力在该日内起伏。微网必须在逐时段满足负荷的前提下，决定每个时段向外网购多少电、储能何时充、何时放，使全天购电费用最小。与后续三问相比，问题一的特殊性在于\textbf{时间尺度最短、信息最完整}：整个优化只覆盖一天，且电价与负荷在制定计划时均已完全已知，因此它是一个标准的确定性逐时段线性规划，不需要任何预测环节。

真正决定题目性质的是两条附加条件。第一，题面要求储能设备在 0:00 与 24:00 的储电量相同，这使储能只能“当天借、当当天还”，无法把电能跨日搬运。第二，附录给出的充放电效率均为 90\%，往返效率仅 81\%，这意味着储能每一轮充放都要付出约 19\% 的循环损耗。两条条件共同把问题一从“无成本套利”变成“有损耗的套利”：只有当充放电时段的电价差足以覆盖循环损耗时，储能才值得投入运行。这一判断将直接决定最优解中储能是否被使用、以及使用到什么程度，是下文机制分析的出发点。

\subsubsection{模型的建立}

以 10 分钟为一时段，全天共 144 个时段。记时段 $t$ 的电价为 $p_t$（元/kWh）、负载功率为 $L_t$、光伏出力为 $P^{pv}_t$（kW），时段长 $\Delta t=1/6$ h。决策变量共五组：

\begin{center}
\small
\begin{tabular}{clc}
\toprule
变量 & 含义 & 单位 \\
\midrule
$b_t$ & 计划购电量 & kWh \\
$c_t$ & 储能充电量（并网点侧） & kWh \\
$q^{dis}_t$ & 储能放电量（并网点侧） & kWh \\
$s_t$ & 弃光量 & kWh \\
$E_t$ & 时段末储电量 & kWh \\
\bottomrule
\end{tabular}
\end{center}

目标为全天购电费用最小：
\begin{equation*}
\min\;\sum_{t=1}^{144} p_t\,b_t .
\end{equation*}

约束由四组构成。其一是逐时段的电量平衡：购电量、光伏电量与放电量之和，等于负载电量、充电量与弃光量之和，
\begin{equation*}
b_t+P^{pv}_t\Delta t+q^{dis}_t=L_t\Delta t+c_t+s_t ,
\end{equation*}
该式为等式，因此“微网提供的电能不低于小区负载”自动成立，同时充电只与弃光竞争、不与负载竞争。其二是储能的状态转移与容量约束，
\begin{equation*}
E_t=E_{t-1}+0.9\,c_t-\frac{q^{dis}_t}{0.9},\qquad 1200\le E_t\le 10800 ,
\end{equation*}
其中 0.9 的两次出现对应充电与放电各自 90\% 的效率。其三是功率约束：附录规定最大充放电功率为 5000 kW，且该限制同时作用于并网点侧与电池侧，两侧均不超过 5000 kW，折算到 10 分钟电量后为 $c_t\le 833.33$ kWh、$q^{dis}_t\le 750.00$ kWh。其四是日周期约束，
\begin{equation}
E_0=E_{144}=6000 ,
\end{equation}
其中 6000 kWh 为附录给定的 1 月 1 日 0:00 初始电量。所有变量非负，弃光量另受物理盈余上界约束。

模型为 720 个变量、288 条等式约束的线性规划，是四问中规模最小的一个，用 HiGHS 在 CPU 上直接求解。

\subsubsection{模型的求解与结果}

最优全天购电费为 $C^{*}=35\,126.948589$ 元，全天购电量 $59\,482.698998$ kWh，当日净负荷（负载电量减光伏电量）为 $55\,541.972433$ kWh；二者之差来自储能循环损耗。求解器返回最优状态（364 次迭代），等式残差不超过 $4.5\times10^{-13}$，变量界零越界。

按题面表 1、表 2 的格式，指定时段的结果列于表~\ref{tab:p1a}，储能在各时段的充放电量与两个端点储电量列于表~\ref{tab:p1b}。六个指定时段的充电量与放电量分别列示、不做冲抵；全天购电费即为上文的 $C^{*}$。

\begin{table}[htbp]
\centering
\small
\caption{问题一：指定时段的购电量与全天合计}
\label{tab:p1a}
\begin{tabular}{lr}
\toprule
时段 & 购电量 / kWh \\
\midrule
10:00--10:10 & 0.00 \\
12:00--12:10 & 486.4029 \\
14:00--14:10 & 0.00 \\
16:00--16:10 & 394.9315 \\
18:00--18:10 & 636.9826 \\
20:00--20:10 & 0.00 \\
\midrule
全天购电量 / kWh & 59\,482.698998 \\
全天购电费 / 元 & 35\,126.948589 \\
\bottomrule
\end{tabular}
\end{table}

\begin{table}[htbp]
\centering
\small
\caption{问题一：储能充放电量与端点储电量}
\label{tab:p1b}
\begin{tabular}{lrr}
\toprule
时段（4 小时块） & 充电量 / kWh & 放电量 / kWh \\
\midrule
0:00--4:00 & 4500.00 & 0.00 \\
4:00--8:00 & 833.33 & 6365.84 \\
8:00--12:00 & 4787.96 & 1703.00 \\
12:00--16:00 & 5286.04 & 91.10 \\
16:00--20:00 & 0.00 & 5780.13 \\
20:00--24:00 & 5333.33 & 2859.87 \\
\midrule
0:00 储电量 / kWh & \multicolumn{2}{r}{6000.00} \\
24:00 储电量 / kWh & \multicolumn{2}{r}{6000.00} \\
\bottomrule
\end{tabular}
\end{table}

\subsubsection{策略形态与机制解释}

最优策略呈现清晰的“低价充电、高价放电”形态（图~\ref{fig:p1a}、图~\ref{fig:p1b}）。充电时段的电价均值为 0.6425 元/kWh，放电时段的电价均值为 1.0845 元/kWh，价差 0.4420 元/kWh；电价与充电量的相关系数为 $-0.446$，与放电量的相关系数为 $+0.592$。按电价分位统计，最便宜的 20\% 时段承担了 55.4\% 的充电量，最贵的 20\% 时段承担了 63.7\% 的放电量。

这一形态是否值得执行，取决于价差能否覆盖循环损耗。往返效率 81\% 要求充放电的价差率超过 $1/0.81-1\approx 23.5\%$ 才有套利空间，而本题实测价差率为 $0.4420/0.6425\approx 68.8\%$，远高于该门槛，因此储能被充分调用：全天充电量与放电量分别为 $20\,740.666132$ kWh 与 $16\,799.939567$ kWh，充电峰值达到功率上限 833.33 kWh/时段，储能轨迹在 0:00 与 24:00 精确闭合于 6000 kWh（图~\ref{fig:p1c}）。当日光伏盈余全部被储能吸收，弃光量为 0；最优解中不存在同时充放电的时段，互补性无需借助二元变量即可满足。

\begin{figure}[htbp]
\centering
\includegraphics[width=0.490\textwidth]{fig03_p1_arbitrage.png}
\caption{净放电量与电价分位}
\label{fig:p1a}
\end{figure}

\begin{figure}[htbp]
\centering
\begin{minipage}[t]{0.485\textwidth}
\centering
\includegraphics[width=\linewidth]{fig04_p1_dispatch.png}
\caption{问题一最优计划购电量与储能充放电安排}
\label{fig:p1b}
\end{minipage}\hfill
\begin{minipage}[t]{0.485\textwidth}
\centering
\includegraphics[width=\linewidth]{fig05_p1_soc.png}
\caption{问题一储电量轨迹}
\label{fig:p1c}
\end{minipage}
\end{figure}

\subsubsection{结果检验}

检验从量级、约束与填报三个层次展开。量级上，费用 $35\,126.95$ 元落在独立构造的区间 $[20\,622.73,\,48\,052.05]$ 元内；其中上界 $48\,052.046591$ 元是“完全不使用储能、按分时电价直接购电”的可行策略费用，最优解相对该策略节省 26.898\%，这一比值同时验证了套利的方向与幅度。约束上，68 项独立断言全部通过：功率平衡与状态转移残差为 0，$E_t\in[1200,10800]$，$E_{144}=E_0$，充放电峰值均不越界。填报上，写出的 \texttt{result1.xlsx} 与解序列逐点一致。

\subsubsection{模型族对照}

为确认上述结论不依赖于某一特定建模形式，本文构造了同题的若干变体并逐一求解，结果列于表~\ref{tab:p1ab}。

\begin{table}[htbp]
\centering
\small
\caption{问题一的模型族与建模选择对照}
\label{tab:p1ab}
\begin{tabular}{p{5.4cm}>{\raggedleft\arraybackslash}p{4.6cm}>{\raggedleft\arraybackslash}p{3.6cm}}
\toprule
对照项 & 费用 / 元 & 相对 $C^{*}$ \\
\midrule
基准线性规划（连续松弛） & 35\,126.948589 & $0.000$\% \\
加 144 个互补二元变量的整数模型 & 35\,126.948589 & $0.000$\% \\
储电量离散化 DP（网格步长 400 kWh） & 37\,616.52 & $+7.087$\% \\
储电量离散化 DP（网格步长 200 kWh） & 36\,203.04 & $+3.063$\% \\
储电量离散化 DP（网格步长 100 kWh） & 35\,688.92 & $+1.600$\% \\
储电量离散化 DP（网格步长 50 kWh） & 35\,404.67 & $+0.791$\% \\
去掉日周期约束 & 32\,909.865256 & $-6.312$\% \\
删除弃光变量 & 35\,126.948589 & $0.000$\% \\
效率作用位置：仅充电侧计效 & 33\,416.53 & $-4.869$\% \\
效率作用位置：仅放电侧计效 & 34\,168.61 & $-2.728$\% \\
效率作用位置：往返整体计效 & 33\,801.50 & $-3.773$\% \\
\bottomrule
\end{tabular}
\end{table}

三点值得强调。其一，加入 144 个互补二元变量的混合整数模型给出的目标值与线性规划\textbf{完全相同}（最优间隙为 0、仅 1 个节点），说明连续松弛在本问上无损，无需付出整数变量的计算代价。其二，储电量离散化的动态规划在四档网格上的目标值恒不低于 $C^{*}$，且随网格加密单调收敛（最细网格偏差 $+0.791\%$），差异来自离散化本身而非模型结构。其三，日周期约束的经济价值为 $2217.083333$ 元，占 $C^{*}$ 的 6.312\%：去掉该约束后模型会把期末库存放空至下限，相当于把日初存量的电能当日变现。这说明“0:00 与 24:00 储电量相同”并非形式要求，而有实质经济含义，也是问题二允许跨日携带电量后费用结构变化的根源。

效率作用位置的三组对照（仅充电计效、仅放电计效、往返整体计效）所得费用比统一处理方式低 958--1710 元，表明本文采用的“充放两侧各按 90\% 计”是几种可能处理方式中最保守的一种。附录给定的 90\% 是题面硬参数；实际电池效率随荷电状态与老化变化，恒效率属工程简化\cite{ref-vykhodtsev-2022-bess-review,ref-huang-2022-dynamic-efficiency}。目标中不计折旧同样是显式取舍——日前微网调度文献常把退化成本纳入目标\cite{ref-zhang-2015-dayahead-degradation}，题面未给相应系数，故本文在假设中声明而不隐藏。

\subsubsection{稳健性}

在 948 个情景（947 个可行）的扰动实验下，问题一的结论判为“稳定”，全部检验标准通过。费用对往返效率最敏感且方向单调：$\eta$ 由 0.9 降至 0.8 时费用上升 8.885\%，升至 1.0 时下降 7.828\%，弹性约为 $-0.70$ 至 $-0.80$；其余参数的影响均不超过 0.7\%，其中初始储电量在 $[3000,9000]$ kWh 内的变化对费用影响不超过 0.07\%——日周期闭合使初值只影响日内可调度裕度，不影响总量。噪声实验显示费用区间半宽不超过 11.6\%，且均值存在负向偏移，该偏移源于价格函数的凹性，不应解读为预测误差可以省钱。
